% OpenStax Calculus Volume 3 -- Section 4.5 Exercises: The Chain Rule
% Exercises reproduced from OpenStax, licensed under CC BY 4.0.
% Source: https://openstax.org/books/calculus-volume-3/pages/4-5-the-chain-rule
\documentclass{exercises}
\title{OpenStax Calculus Volume 3}
\subtitle{Section 4.5 Exercises: The Chain Rule}
\sourceurl{https://openstax.org/books/calculus-volume-3/pages/4-5-the-chain-rule}
\author{}
\date{}

\begin{document}
\maketitle


For the following exercises, use the information provided to solve the problem.

\begin{enumerate}[start=1]
\item Let $w(x,y,z)=xy\cos z$, where $x=t,y=t^{2}$, and $z=\arcsin t$. Find $\frac{dw}{dt}$.
\item Let $w(t,v)=e^{tv}$ where $t=r+s$ and $v=rs$. Find $\frac{\partial w}{\partial r}$ and $\frac{\partial w}{\partial s}$.
\item If $w=5x^{2}+2y^{2},x=-3s+t$, and $y=s-4t$, find $\frac{\partial w}{\partial s}$ and $\frac{\partial w}{\partial t}$.
\item If $w=xy^{2},x=5\cos(2t)$, and $y=5\sin(2t)$, find $\frac{dw}{dt}$.
\item If $f(x,y)=xy,x=r\cos \theta$, and $y=r\sin \theta$, find $\frac{\partial f}{\partial r}$ and express the answer in terms of $r$ and $\theta$.
\item Suppose $f(x,y)=x+y$, where $x=r\cos \theta$ and $y=r\sin \theta$. Find $\frac{\partial f}{\partial \theta}$.
\end{enumerate}

For the following exercises, find $\frac{df}{dt}$ using the chain rule and direct substitution.

\begin{enumerate}[resume]
\item $f(x,y)=x^{2}+y^{2}$, $x=t,y=t^{2}$
\item $f(x,y)=\sqrt{x^{2}+y^{2}},y=t^{2},x=t$
\item $f(x,y)=xy,x=1-\sqrt{t},y=1+\sqrt{t}$
\item $f(x,y)=\frac{x}{y},x=e^{t},y=2e^{t}$
\item $f(x,y)=\ln(x+y)$, $x=e^{t},y=e^{t}$
\item $f(x,y)=x^{4}$, $x=t,y=t$
\item Let $w(x,y,z)=x^{2}+y^{2}+z^{2}$, $x=\cos t,y=\sin t$, and $z=e^{t}$. Express $w$ as a function of $t$ and find $\frac{dw}{dt}$ directly. Then, find $\frac{dw}{dt}$ using the chain rule.
\item Let $z=x^{2}y$, where $x=t^{2}$ and $y=t^{3}$. Find $\frac{dz}{dt}$.
\item Let $u=e^{x}\sin y$, where $x=-\ln(2t)$ and $y=\pi t$. Find $\frac{du}{dt}$ when $x=\ln 2$ and $y=\frac{\pi}{4}$.
\end{enumerate}

For the following exercises, find $\frac{dy}{dx}$ using partial derivatives.

\begin{enumerate}[resume]
\item $\sin(6x)+\tan(8y)+5=0$
\item $x^{3}+y^{2}x-3=0$
\item $\sin(x+y)+\cos(x-y)=4$
\item $x^{2}-2xy+y^{4}=4$
\item $xe^{y}+ye^{x}-2x^{2}y=0$
\item $x^{2/3}+y^{2/3}=a^{2/3}$
\item $x\cos(xy)+y\cos x=2$
\item $e^{xy}+ye^{y}=1$
\item $x^{2}y^{3}+\cos y=0$
\item Find $\frac{dz}{dt}$ using the chain rule where $z=3x^{2}y^{3},x=t^{4}$, and $y=t^{2}$.
\item Let $z=3\cos x-\sin(xy),x=\frac{1}{t}$, and $y=3t$. Find $\frac{dz}{dt}$.
\item Let $z=e^{1-xy},x=t^{1/3}$, and $y=t^{3}$. Find $\frac{dz}{dt}$.
\item Find $\frac{dz}{dt}$ by the chain rule where $z=\cosh^{2}(xy),x=\frac{1}{2}t$, and $y=e^{t}$.
\item Let $z=\frac{x}{y},x=2\cos u$, and $y=3\sin v$. Find $\frac{\partial z}{\partial u}$ and $\frac{\partial z}{\partial v}$.
\item Let $z=e^{x^{2}y}$, where $x=\sqrt{uv}$ and $y=\frac{1}{v}$. Find $\frac{\partial z}{\partial u}$ and $\frac{\partial z}{\partial v}$.
\item If $z=xye^{x/y}$, $x=r\cos \theta$, and $y=r\sin \theta$, find $\frac{\partial z}{\partial r}$ and $\frac{\partial z}{\partial \theta}$ when $r=2$ and $\theta=\frac{\pi}{6}$.
\item Find $\frac{\partial w}{\partial s}$ if $w=4x+y^{2}+z^{3},x=e^{rs^{2}},y=\ln(\frac{r+s}{t})$, and $z=rst^{2}$.
\item If $w=\sin(xyz),x=1-3t,y=e^{1-t}$, and $z=4t$, find $\frac{\partial w}{\partial t}$.
\end{enumerate}

For the following exercises, use this information: A function $f(x,y)$ is said to be homogeneous of degree $n$ if $f(tx,ty)=t^{n}f(x,y)$. For all homogeneous functions of degree $n$, the following equation is true: $x\frac{\partial f}{\partial x}+y\frac{\partial f}{\partial y}=nf(x,y)$. Show that the given function is homogeneous and verify that $x\frac{\partial f}{\partial x}+y\frac{\partial f}{\partial y}=nf(x,y)$.

\begin{enumerate}[resume]
\item $f(x,y)=3x^{2}+y^{2}$
\item $f(x,y)=\sqrt{x^{2}+y^{2}}$
\item $f(x,y)=x^{2}y-2y^{3}$
\item The volume of a right circular cylinder is given by $V(x,y)=\pi x^{2}y$, where $x$ is the radius of the cylinder and $y$ is the cylinder height. Suppose $x$ and $y$ are functions of $t$ given by $x=\frac{1}{2}t$ and $y=\frac{1}{3}t$ so that $x$ and $y$ are both increasing with time. How fast is the volume increasing when $x=2$ and $y=\frac{4}{3}$?
\item The pressure $P$ of a gas is related to the volume and temperature by the formula $PV=kT$, where temperature is expressed in kelvins. Express the pressure of the gas as a function of both $V$ and $T$. Find $\frac{dP}{dt}$ when $k=1$, $\frac{dV}{dt}=2$ cm$^{3}$/min, $\frac{dT}{dt}=\frac{1}{2}$ K/min, $V=20$ cm$^{3}$, and $T=20^\circ \text{F}$.
\item The radius of a right circular cone is increasing at $3$ cm/min whereas the height of the cone is decreasing at $2$ cm/min. Find the rate of change of the volume of the cone when the radius is $13$ cm and the height is $18$ cm.
\item The volume of a frustum of a cone is given by the formula $V=\frac{1}{3}\pi z(x^{2}+y^{2}+xy)$, where $x$ is the radius of the smaller circle, $y$ is the radius of the larger circle, and $z$ is the height of the frustum (see figure). Find the rate of change of the volume of this frustum when $x=10$ in., $y=12$ in., and $z=18$ in. if $\frac{dz}{dt}=-5,\frac{dx}{dt}=1,\frac{dy}{dt}=1$ (all in/min). \\ \textit{[Figure: A conical frustum (that is, a cone with the pointy end cut off) with height $z$, larger radius $y$, and smaller radius $x$.]}
\item A closed box is in the shape of a rectangular solid with dimensions $x$, $y$, and $z$. (Dimensions are in inches.) Suppose each dimension is changing at the rate of $0.5$ in./min. Find the rate of change of the total surface area of the box when $x=2$ in., $y=3$ in., and $z=1$ in.
\item The total resistance in a circuit that has three individual resistances represented by $x,y$, and $z$ is given by the formula $R(x,y,z)=\frac{xyz}{yz+xz+xy}$. Suppose at a given time the $x$ resistance is $100\Omega$, the $y$ resistance is $200\Omega$, and the $z$ resistance is $300\Omega$. Also, suppose the $x$ resistance is changing at a rate of $2\Omega/\text{min}$, the $y$ resistance is changing at the rate of $1\Omega/\text{min}$, and the $z$ resistance has no change. Find the rate of change of the total resistance in this circuit at this time.
\item The temperature $T$ at a point $(x,y)$ is $T(x,y)$ and is measured using the Celsius scale. A fly crawls so that its position after $t$ seconds is given by $x=\sqrt{1+t}$ and $y=2+\frac{1}{3}t$, where $x$ and $y$ are measured in centimeters. The temperature function satisfies $T_{x}(2,3)=4$ and $T_{y}(2,3)=3$. How fast is the temperature increasing on the fly's path after $3$ sec?
\item The $x$ and $y$ components of a fluid moving in two dimensions are given by the following functions: $u(x,y)=2y$ and $v(x,y)=-2x$; $x\ge0$, $y\ge0$. The speed of the fluid at the point $(x,y)$ is $s(x,y)=\sqrt{u(x,y)^{2}+v(x,y)^{2}}$. Find $\frac{\partial s}{\partial x}$ and $\frac{\partial s}{\partial y}$ using the chain rule.
\item Let $u=u(x,y,z)$, where $x=x(w,t),y=y(w,t),z=z(w,t),w=w(r,s)$, and $t=t(r,s)$. Use a tree diagram and the chain rule to find an expression for $\frac{\partial u}{\partial r}$.
\end{enumerate}

\end{document}
